\chapter{微分と積分の関係}
    \section{\texorpdfstring{$\lim_{x\to\infty}{f(x)}$}{lim x→∞ f(x)}が存在すれば\texorpdfstring{$\int_{0}^{\infty}{f'(x) dx} = f(\infty)-f(0)$}{∫\_0\string^∞ f'(x) = f(∞) - f(0)}}
        \begin{shadebox}
            $\textup{C}^1$級関数$f(x)$について$\displaystyle \lim_{x\to\infty}{f(x)}$が存在するとき$f'(x) $は$0\sim\infty$で可積分である。すなわち$\displaystyle \int_{0}^{\infty}{f'(x) dx}$が存在する。
        \end{shadebox}
        \begin{proof}
            \quad\par
            $\displaystyle f(\infty) \coloneqq \lim_{x\to\infty}{f(x)}$とすると任意の$\epsilon>0$に対してある正数$X(\epsilon)$が存在して$x \geq X \Rightarrow |f(x)-f(\infty)| < \epsilon$であるから、$x\geq X$とすると
            \[|[f(x)-f(0)]-[f(\infty)-f(0)]| < \epsilon\]
            \[\therefore \; \left|\int_{0}^{x}{f'(u) du}-[f(\infty)-f(0)]\right| < \epsilon\]
            $\epsilon \to 0$のとき$X(\epsilon) \to \infty$であるから結局
            \[\int_{0}^{\infty}{f'(x) dx} = f(\infty)-f(0)\]
        \end{proof}
    \section{\texorpdfstring{$\derivLong{\integrate{a}{x}{f(x)}{}{x}}{x}{} = f(x)$}{(d/dx)∫\_a\string^x f(x) dx = f(x)}}
        \begin{shadebox}
            連続関数$f(x)$について
            \[\derivLong{\integrate{a}{x}{f(x)}{}{x}}{x}{} = f(x)\]
        \end{shadebox}
        \begin{proof}
            \begin{equation}
                \derivLong{\integrate{a}{x}{f(x)}{}{x}}{x}{}= \lim_{h\rightarrow0}{\frac{1}{h}\left(\integrate{a}{x+h}{f(x)}{}{x}-\integrate{a}{x}{f(x)}{}{x}\right)} =  \lim_{h\rightarrow0}{\frac{1}{h}\integrate{x}{x+h}{f(x)}{}{x}}
            \end{equation}
            $h>0$の場合を示す。$h<0$の場合も同じ要領で示せる。
            積分の平均値の定理よりある$\xi\in(x,x+h)$が存在して
            \[hf(\xi)=\integrate{x}{x+h}{f(x)}{}{x}\]
            となるから
            \[(1) = \lim_{h\rightarrow0}{f(\xi)}=f(x)\]
        \end{proof}
    \section{
        畳み込みの微分\texorpdfstring{\\}{}
        \texorpdfstring{$\derivLong{\integrate{s=0}{x}{f(s)g(x-s)}{}{s}}{x}{} = f(x)g(0) + \integrate{s=0}{x}{f(s)\left.\deriv{g(x)}{x}{}\right|_{x=x-s}}{}{s}$}{}
    }
        \begin{proof}
            \newcommand{\dx}{{\Delta x}}
            \renewcommand{\d}[1]{{\mathrm{d}#1}}
            \quad\par
            $I(x) \coloneqq \integrate{0}{x}{f(s)g(x-s)}{}{s}$とおき、微分の定義に従い$\lim_{\dx\to 0}\frac{I(x+\dx)-I(x)}{\dx}$を計算する。
            \begin{align*}
                I(x+\dx)-I(x) &= \left(\int_{0}^{x}+\int_{x}^{x+\dx}\right)f(s)g(x+\dx-s)\d{s} - \integrate{s=0}{x}{f(s)g(x-s)}{}{s} \\
                &= \integrate{0}{x}{f(s)\left(g(x-s+\dx)-g(x-s)\right)}{}{s} + \integrate{x}{x+\dx}{f(s)g(x+\dx-s)}{}{s}
            \end{align*}
            積分の平均値の定理より$^\exists\theta\in(x,x+\dx)\st\integrate{x}{x+\dx}{f(s)g(x+\dx-s)}{}{s} = \dx f(\theta)g(x+\dx-\theta)$であるから
            \begin{align*}
                \lim_{\dx\to 0}\frac{I(x+\dx)-I(x)}{\dx} &= \lim_{\dx\to 0}\left[f(\theta)g(x+\dx-\theta) + \integrate{0}{x}{f(s)\frac{g(x-s+\dx)-g(x-s)}{\dx}}{}{s}\right] \\
                &= f(x)g(0) + \integrate{s=0}{x}{f(s)\left.\deriv{g(x)}{x}{}\right|_{x=x-s}}{}{s}
            \end{align*}
        \end{proof}
    \section{
        重積分の微分\texorpdfstring{\\}{}
        \small{
            \texorpdfstring{$\derivLong{\left(\integrate{x_1=0}{x}{\dotsi\integrate{x_n=0}{x-x_1-\dotsb-x_{n-1}}{f_1(x_1)\dotsm f_n(x_n)}{}{x_n}\dotsi}{}{x_1}\right)}{x}{}$}{} \texorpdfstring{\\}{}
            \texorpdfstring{$= \int_{x_1=0}^x\dotsi\int_{x_{n-1}=0}^{x-x_1-\dotsb -x_{n-2}} f_1(x_1)\dotsm f_{n-1}(x_{n-1})f_n(x-x_1-\dotsb-x_{n-1}) \mathrm{d}x_{n-1}\dotsi\mathrm{d}x_{1}$}{}
        }
    }
        \label{重積分の微分}
        \begin{shadebox}
            $\omega>0$とする。
            \[\Omega\coloneqq\{(x_1,\dotsc,x_n)|x_1,\dotsc,x_n\geq 0, x_1+\dotsb+x_n \leq \omega\}\]
            \[\overline{\Omega}^+\coloneqq\{(x_1,\dotsc,x_n)|x_1,\dotsc,x_n\geq 0, x_1+\dotsb+x_n = \omega\} \subset \Omega\]
            $f_1(x),\dotsc,f_n(x)$を$\{x|x\geq 0\}$上の連続関数とし、$f(\bm{x})\coloneqq f_1(x_1)\dotsm f_n(x_n),\quad F(\omega) \coloneqq \integrate{\Omega}{}{f(\bm{x})}{}{\bm{x}}$とするとき、次が成り立つ。
            \[\deriv{F(\omega)}{\omega}{} = \integrate{\overline{\Omega}^+}{}{f(\bm{x})}{}{\bm{x}}\]
            要するに、領域$x_1,\dotsc,x_n\geq 0$の超平面$x_1+\dotsb+x_n=\omega$の原点側の部分における$f_1(x_1)\dotsm f_n(x_n)$の積分を$\omega$で微分したものは領域$x_1,\dotsc,x_n\geq 0$と超平面の共通部分上での積分に等しい。\\
            \\
            \quad この定理は2つの独立な非負の確率変数の和の確率密度関数を計算する時に役立つ。
        \end{shadebox}
        \begin{proof}
            \quad\par
            \newcommand{\domega}{{\Delta\omega}}
            \renewcommand{\d}[1]{{\mathrm{d}#1}}
            微分の定義に従い$\lim_{\domega \to 0}\frac{F(\omega+\domega)-F(\omega)}{\domega}$を計算する。
            \begin{align*}
                F(\omega+\domega) &= \integrate{x_1=0}{\omega}{\dotsi\integrate{x_n=0}{\omega-x_1-\dotsb-x_{n-1}}{f_1(x_1)\dotsm f_n(x_n)}{}{x_n}\dotsi}{}{x_1} \\
                &= \left(\int_{x_1=0}^\omega+\int_{x_1=\omega}^{\omega+\domega}\right)f_1(x_1)\textcolor{red}{\left(\int_{x_2=0}^{\omega-x_1}+\int_{x_2=\omega-x_1}^{\omega-x_1+\domega}\right)f_2(x_2)\dotsi}\\
                & \textcolor{red}{\left(\int_{x_n=0}^{\omega-x_1-\dotsb-x_{n-1}}+\int_{x_n=\omega-x_1-\dotsb-x_{n-1}}^{\omega-x_1-\dotsb-x_{n-1}+\domega}\right)f_n(x_n)\d{x_n}\dotsi\d{x_2}}\d{x_1}
            \end{align*}
            であり、括弧の中身を全部展開すれば$F(\omega)$を含む$2^n$個の項が出る。
            $\frac{F(\omega+\domega)-F(\omega)}{\domega}$を考えるときは$F(\omega)$は相殺するから放っておけば良い。
            $F(\omega)$を除いた$2^n-1$項を$\domega$で除して$\domega \to 0$としたとき、ある1項を除いて他は0になり、残った項が$\integrate{\overline{\Omega}^+}{}{f(\bm{x})}{}{\bm{x}}$になることを以下で示していく。
            \[g_1(\omega,x_1) \coloneqq \textcolor{red}{\int_{x_2=0}^{\omega-x_1+\domega}f_2(x_2)\dotsi\int_{x_n=0}^{\omega-x_1-\dotsb-x_{n-1}+\domega}f_n(x_n)\d{x_n}\dotsi\d{x_2}}\]
            とおけば
            \begin{align*}
                F(\omega+\domega) = \underbrace{\int_{x_1=0}^{\omega}f_1(x_1)g_1(\omega,x_1)\d{x_1}}_{=:h_{11}(\omega,\domega)} + \underbrace{\int_{x_1=\omega}^{\omega+\domega}f_1(x_1)g_1(\omega,x_1)\d{x_1}}_{=:h_{12}(\omega,\domega)}
            \end{align*}
            となる。$\lim_{\domega\to 0}\frac{h_{12}(\omega,\domega)}{\domega}=0$を示す。
            積分の平均値の定理から$^\exists\theta_1 \in (\omega,\omega+\domega)\st h_{12}(\omega,\domega)=\domega f_1(\theta_1)g_1(\omega,\theta_1)$であるので
            \[\lim_{\domega\to 0}\frac{h_{12}(\omega,\domega)}{\domega} = \lim_{\domega\to 0}f_1(\theta_1)g_1(\omega,\theta_1) = f_1(\omega)\lim_{\domega\to 0}g_1(\omega,\theta_1) = 0\]
            結局$h_{11}(\omega,\domega)$だけ考えれば良い。
            \[g_2(\omega,x_1,x_2) \coloneqq \int_{x_3=0}^{\omega-x_1-x_2+\domega}f_3(x_3)\dotsi\int_{x_n=0}^{\omega-x_1-\dotsb-x_{n-1}+\domega}f_n(x_n)\d{x_n}\dotsi\d{x_3}\]
            とおけば
            \begin{align*}
                h_{11}(\omega,\domega) &= \underbrace{\int_{x_1=0}^\omega f_1(x_1)\int_{x_2=0}^{\omega-x_1}f_2(x_2)g_2(\omega,x_1,x_2)\d{x_2}\d{x_1}}_{=:h_{21}(\omega,\domega)} \\
                &+ \underbrace{\int_{x_1=0}^\omega f_1(x_1)\int_{x_2=\omega-x_1}^{\omega-x_1+\domega}f_2(x_2)g_2(\omega,x_1,x_2)\d{x_2}\d{x_1}}_{=:h_{22}(\omega,\domega)}
            \end{align*}
            となる。$\lim_{\domega\to 0}\frac{h_{22}(\omega,\domega)}{\domega} = 0$を示す。
            \par
            積分の平均値の定理から$^\exists\theta_2\in(\omega-x_1,\omega-x_1+\domega)\st\int_{x_2=\omega-x_1}^{\omega-x_1+\domega}f_2(x_2)g_2(\omega,x_1,x_2)\d{x_2} = \domega f_2(\theta_2)g_2(\omega,x_1,\theta_2)$であるので
            \[\lim_{\domega\to 0}\frac{h_{22}(\omega,\domega)}{\domega} = \lim_{\domega\to 0}\int_{x_1=0}^\omega f_1(x_1)f_2(\theta_2)g_2(\omega,x_1,\theta_2)\d{x_1}\]
            これと$\lim_{\domega\to 0}g_2(\omega,x_1,\theta_2)=0$であり、かつ$\int_{x_1=0}^\omega f_1(x_1)f_2(\theta_2)\d{x_1}$が有界であることから
            \[\lim_{\domega\to 0}\frac{h_{22}(\omega,\domega)}{\domega} = 0\]
            となるので、結局$h_{21}(\omega,\domega)$だけを考えれば良い。
            同様の議論を繰り返すと、結局$h_{ij}(\omega,\domega)\quad((i,j)\in\{1,\dotsc,n-1\}\times\{1,2\})$に対して$\lim_{\domega\to 0}\frac{h_{i2}(\omega,\domega)}{\domega}=0$であり、残る項は
            \begin{align*}
                &\quad h_{n-1,1}(\omega,\domega) \\
                &= \int_{x_1=0}^{\omega}f_1(x_1)\dotsi\int_{x_{n-1}=0}^{\omega-x_1-\dotsb-x_{n-2}}f_{n-1}(x_{n-1}) \\
                & \left(\int_{x_n=0}^{\omega-x_1-\dotsb-x_{n-1}}+\int_{x_n=\omega-x_1-\dotsb-x_{n-1}}^{\omega-x_1-\dotsb-x_{n-1}+\domega}\right)f_n(x_n)\d{x_n}\d{x_{n-1}}\dotsi\d{x_1} \\
                &= F(\omega) + \underbrace{\int_{x_1=0}^{\omega}f_1(x_1)\dotsi\int_{x_{n-1}=0}^{\omega-x_1-\dotsb-x_{n-2}}f_{n-1}(x_{n-1})\textcolor{red}{\int_{x_n=\omega-x_1-\dotsb-x_{n-1}}^{\omega-x_1-\dotsb-x_{n-1}+\domega}f_n(x_n)\d{x_n}}\d{x_{n-1}}\dotsi\d{x_1}}_{=:h_{n2}(\omega,\domega)}
            \end{align*}
            だけである。
            \par
            積分の平均値の定理から$^\exists\theta_n\in(\omega-x_1-\dotsb-x_{n-1},\omega-x_1-\dotsb-x_{n-1}+\domega)\st\textcolor{red}{\int_{x_n=\omega-x_1-\dotsb-x_{n-1}}^{\omega-x_1-\dotsb-x_{n-1}+\domega}f_n(x_n)\d{x_n}} = \domega f_n(\theta_n)$であるので
            \begin{align*}
                \lim_{\domega\to 0}\frac{h_{n2}(\omega,\domega)}{\domega} &= \int_{x_1=0}^{\omega}f_1(x_1)\dotsi\int_{x_{n-1}=0}^{\omega-x_1-\dotsb-x_{n-2}}f_{n-1}(x_{n-1})f_n(\theta_n)\d{x_{n-1}}\dotsi\d{x_1} \\
                &= \int_{x_1=0}^{\omega}f_1(x_1)\dotsi\int_{x_{n-1}=0}^{\omega-x_1-\dotsb-x_{n-2}}f_{n-1}(x_{n-1})f_n(\omega-x_1-\dotsb-x_{n-1})\d{x_{n-1}}\dotsi\d{x_1} \\
                &= \integrate{\overline{\Omega}^+}{}{f(\bm{x})}{}{\bm{x}}
            \end{align*}
            以上より
            \[\deriv{F(\omega)}{\omega}{} = \lim_{\domega \to 0}\frac{F(\omega+\domega)-F(\omega)}{\domega} = \lim_{\domega\to 0}\frac{h_{n2}(\omega,\domega)}{\domega} = \integrate{\overline{\Omega}^+}{}{f(\bm{x})}{}{\bm{x}}\]
        \end{proof}
    \section{
        系\texorpdfstring{\\}{}
        \small{
            \texorpdfstring{$\derivLong{\left(\integrate{x_1=-\infty}{\infty}{\dotsi\integrate{x_n=-\infty}{x-x_1-\dotsb-x_{n-1}}{f_1(x_1)\dotsm f_n(x_n)}{}{x_n}\dotsi}{}{x_1}\right)}{x}{}$\\}{}
            \texorpdfstring{$= \integrate{x_1=-\infty}{\infty}{\dotsi\integrate{x_{n-1}=-\infty}{x-x_1-\dotsb-x_{n-2}}{f_1(x_1)\dotsm f_{n-1}(x_{n-1})f_n(x-x_1-\dotsb-x_{n-1})}{}{x_{n-1}}\dotsi}{}{x_1}$}{}
        }
    }
        \label{重積分の微分の系}
        \begin{shadebox}
            \[\Omega\coloneqq\{(x_1,\dotsc,x_n)|x_1+\dotsb+x_n \leq \omega\}\]
            \[\overline{\Omega}^+\coloneqq\{(x_1,\dotsc,x_n)|x_1+\dotsb+x_n = \omega\} \subset \Omega\]
            $f_1(x),\dotsc,f_n(x)$を$\mathbb{R}$上の連続関数とし、$f(\bm{x})\coloneqq f_1(x_1)\dotsm f_n(x_n),\quad F(\omega) \coloneqq \integrate{\Omega}{}{f(\bm{x})}{}{\bm{x}}$とするとき、次が成り立つ。
            \[\deriv{F(\omega)}{\omega}{} = \integrate{\overline{\Omega}^+}{}{f(\bm{x})}{}{\bm{x}}\]
            要するに、超平面$x_1+\dotsb+x_n=\omega$の負領域における$f_1(x_1)\dotsm f_n(x_n)$の積分を$\omega$で微分したものは超平面上での積分に等しい。\\
            \\
            \quad この定理は2つの独立な確率変数の和の確率密度関数を計算する時に役立つ。
        \end{shadebox}
        \begin{proof}
            \quad\par
            直前の定理の証明を少し改造するだけで簡単に証明できる。
        \end{proof}